\documentclass{article}
\usepackage[utf8]{inputenc}
\usepackage{geometry}
\usepackage{enumitem}
\usepackage{hyperref}
\usepackage{xcolor}
\usepackage{graphicx}
\usepackage{array}
\usepackage{multirow}
\usepackage{amsmath, amssymb}
\usepackage{chemfig}
\usepackage{mhchem}
\usepackage{tikz}
\usepackage{setspace}
\usepackage{soul}
\usepackage{mathtools}
\usepackage{booktabs}
\usepackage{chemformula}
\usepackage{siunitx}
\usetikzlibrary{positioning, calc}
\usepackage{longtable}
\geometry{margin=1in}
\setlength{\parskip}{0.5em}
\usepackage{cancel}


\begin{document}

\begin{center}
    \vspace*{-1cm}
    \begin{tabular}{p{12cm} r}
        & \textbf{Name:} \underline{\hspace{5cm}} \\
    \end{tabular}

    \vspace{0.2cm}
    {\LARGE \textbf{Final Review Worksheet}}
\end{center}    

\noindent\begin{tabular}{|p{0.9\linewidth}|}
\hline
\textbf{(1)} For the reaction $\text{A} + \text{heat} \rightarrow \text{B}$, which statement is correct? \\[0.2cm]
\begin{minipage}[t]{0.85\linewidth}
\begin{enumerate}[label=\alph*.]
\item Adding heat increases [B]
\item Adding heat decreases [B]
\item Removing heat increases [B]
\item Heat has no effect on equilibrium
\end{enumerate}
\end{minipage}\\[0.2cm]
\hline
\end{tabular}

\noindent\begin{tabular}{|p{0.9\linewidth}|}
\hline
\textbf{(2)} Which of the following best describes a saturated solution at equilibrium?\\[0.2cm]
\begin{minipage}[t]{0.85\linewidth}
\begin{enumerate}[label=\alph*.]
\item Solute is dissolving faster than it crystallizes
\item Solute and solvent are reacting chemically
\item Dissolving and crystallization occur at equal rates
\item No solute particles are moving
\end{enumerate}
\end{minipage}\\[0.2cm]
\hline
\end{tabular}

\noindent\begin{tabular}{|p{0.9\linewidth}|}
\hline
\textbf{(3)} Calcium hydroxide reacts with hydrochloric acid to form calcium chloride and water. Which statement best describes the acid-base behavior occurring in this reaction?\\[0.2cm]
Ca(OH)$_2$ + 2HCl $\rightarrow$ CaCl$_2$ + 2H$_2$O\\[0.2cm]
\begin{minipage}[t]{0.85\linewidth}
\begin{enumerate}[label=\alph*.]
\item HCl donates protons to OH$^-$ ions, forming water as Ca$^{2+}$ pairs with Cl$^-$
\item Ca(OH)$_2$ donates protons to HCl, forming hydrogen gas
\item Both reactants accept protons, preventing water formation
\item Ca$^{2+}$ acts as the proton donor while Cl$^-$ acts as the proton acceptor
\end{enumerate}
\end{minipage}\\[0.2cm]
\hline
\end{tabular}

\noindent\begin{tabular}{|p{0.9\linewidth}|}
\hline
\textbf{(4)} When comparing HF and HCl, HF has a much higher boiling point. Which force best explains this difference?\\[0.2cm]
\begin{minipage}[t]{0.85\linewidth}
\begin{enumerate}[label=\alph*.]
\item Dipole-dipole interactions
\item Hydrogen bonding
\item London dispersion forces
\item Covalent bonding
\end{enumerate}
\end{minipage}\\[0.2cm]
\hline
\end{tabular}

\noindent\begin{tabular}{|p{0.9\linewidth}|}
\hline
\textbf{(5)} Which of the following increases when temperature increases for a gas at constant volume?\\[0.2cm]

\begin{minipage}[t]{0.85\linewidth}
\begin{enumerate}[label=\alph*.]
\item Pressure
\item Molecular mass
\item Intermolecular forces
\item Volume
\end{enumerate}
\end{minipage}\\[0.2cm]
\hline
\end{tabular}

\noindent\begin{tabular}{|p{0.9\linewidth}|}
\hline
\textbf{(6)} A chemist has 50.0~mL of 2.0~M HCl solution. She adds water to dilute it to a final volume of 200.0~mL. What is the molarity of the diluted solution?\\[0.2cm]
\begin{minipage}[t]{0.85\linewidth}
\begin{enumerate}[label=\alph*.]
\item 0.25~M
\item 0.50~M
\item 1.0~M
\item 2.0~M
\end{enumerate}
\end{minipage}\\[0.2cm]
\hline
\end{tabular}

\noindent\begin{tabular}{|p{0.9\linewidth}|}
\hline
\textbf{(7)} An unsaturated solution contains less solute than the maximum that can dissolve at that temperature.\\[0.2cm]
\textbf{True} \hspace{1cm} \textbf{False}\\[0.2cm]
\hline
\end{tabular}

\noindent\begin{tabular}{|p{0.9\linewidth}|}
\hline
\textbf{(8)} What happens to an atom's oxidation number when it is oxidized?\\[0.2cm]
\begin{minipage}[t]{0.85\linewidth}
\begin{enumerate}[label=\alph*.]
\item It increases
\item It decreases
\item It stays the same
\item It becomes zero
\end{enumerate}
\end{minipage}\\[0.2cm]
\hline
\end{tabular}

\noindent\begin{tabular}{|p{0.9\linewidth}|}
\hline
\textbf{(9)} Which of the following solutes is generally soluble in water according to solubility rules?\\[0.2cm]
\begin{minipage}[t]{0.85\linewidth}
\begin{enumerate}[label=\alph*.]
\item BaSO$_4$
\item AgCl
\item KNO$_3$
\item PbI$_2$
\end{enumerate}
\end{minipage}\\[0.2cm]
\hline
\end{tabular}

\noindent\begin{tabular}{|p{0.9\linewidth}|}
\hline
\textbf{(10)} If $Q > K_{eq}$ for the reaction $\text{F} \rightleftharpoons \text{G}$, which direction will the system shift?\\[0.2cm]
\begin{minipage}[t]{0.85\linewidth}
\begin{enumerate}[label=\alph*.]
\item Toward reactants
\item Toward products
\item No shift
\item The reaction stops
\end{enumerate}
\end{minipage}\\[0.2cm]
\hline
\end{tabular}

\noindent\begin{tabular}{|p{0.9\linewidth}|}
\hline
\textbf{(11)} 0.250~mol of oxygen gas is collected at 740.~mmHg and 18.0$^\circ$C in a 2.00~L flask. What would be the pressure at 0.500~mol in same flask and same T? (assume ideal)\\[0.2cm]

\begin{minipage}[t]{0.85\linewidth}
\begin{enumerate}[label=\alph*.]
\item 1480~mmHg
\item 740.~mmHg
\item 370.~mmHg
\item 1110~mmHg
\end{enumerate}
\end{minipage}\\[0.2cm]
\hline
\end{tabular}

\noindent\begin{tabular}{|p{0.9\linewidth}|}
\hline
\textbf{(12)} A Brönsted-Lowry acid is defined as a substance that: \\[0.2cm]
\begin{minipage}[t]{0.85\linewidth}
\begin{enumerate}[label=\alph*.]
\item Donates a proton
\item Accepts a proton
\item Releases electrons
\item Accepts electrons
\end{enumerate}
\end{minipage}\\[0.2cm]
\hline
\end{tabular}

\noindent\begin{tabular}{|p{0.9\linewidth}|}
\hline
\textbf{(13)} A buffer solution is prepared using 0.120~mol NH$_3$ and 0.080~mol NH$_4^+$ (p$K_b=4.75$) in 1.00~L. After calculating the pH, what is the corresponding $[\mathrm{OH^-}]$?\\[0.2cm]
\begin{minipage}[t]{0.85\linewidth}
\begin{enumerate}[label=\alph*.]
\item pH = 9.32,\ $[\mathrm{OH^-}] = 2.1\times10^{-5}$~M
\item pH = 9.03,\ $[\mathrm{OH^-}] = 1.1\times10^{-5}$~M
\item pH = 9.48,\ $[\mathrm{OH^-}] = 3.0\times10^{-5}$~M
\item pH = 8.90,\ $[\mathrm{OH^-}] = 7.9\times10^{-6}$~M
\end{enumerate}
\end{minipage}\\[0.2cm]
\hline
\end{tabular}

\noindent\begin{tabular}{|p{0.9\linewidth}|}
\hline
\textbf{(14)} Which of the following would double the (m/v)\% of a solution?\\[0.2cm]
\begin{minipage}[t]{0.85\linewidth}
\begin{enumerate}[label=\alph*.]
\item Double solute, same volume
\item Double volume, same solute
\item Halve solute, double volume
\item Keep solute and volume constant
\end{enumerate}
\end{minipage}\\[0.2cm]
\hline
\end{tabular}

\noindent\begin{tabular}{|p{0.9\linewidth}|}
\hline
\textbf{(15)} A 5\% (m/v) glucose solution contains 5~g of glucose in:\\[0.2cm]
\begin{minipage}[t]{0.85\linewidth}
\begin{enumerate}[label=\alph*.]
\item 5~mL of water
\item 100~mL of water
\item 100~mL of solution
\item 95~mL of solution
\end{enumerate}
\end{minipage}\\[0.2cm]
\hline
\end{tabular}

\noindent\begin{tabular}{|p{0.9\linewidth}|}
\hline
\textbf{(16)} The density of aluminum is $2.70~\text{g/cm}^3$. What is the mass of a block of aluminum with a volume of $15.0~\text{cm}^3$? \\
\begin{minipage}[t]{0.85\linewidth}
\begin{enumerate}[label=\alph*.]
    \item $40.5~\text{g}$
    \item $5.56~\text{g}$
    \item $17.7~\text{g}$
    \item $2.70~\text{g}$
\end{enumerate}
\end{minipage}\\[0.2cm]
\hline
\end{tabular}

\noindent\begin{tabular}{|p{0.9\linewidth}|}
\hline
\textbf{(17)} Convert 0.750 L into cubic centimeters. (Use $1~\text{L} = 1000~\text{cm}^3$) \\
\begin{minipage}[t]{0.85\linewidth}
\begin{enumerate}[label=\alph*.]
    \item $75.0~\text{cm}^3$
    \item $750.~\text{cm}^3$
    \item $7.50~\text{cm}^3$
    \item $0.00750~\text{cm}^3$
\end{enumerate}
\end{minipage}\\[0.2cm]
\hline
\end{tabular}

\noindent\begin{tabular}{|p{0.9\linewidth}|}
\hline
\textbf{(18)} Identify the reaction type:\\[0.2cm]
2H$_2$O$_2$ $\rightarrow$ 2H$_2$O + O$_2$\\[0.2cm]
\begin{minipage}[t]{0.85\linewidth}
\begin{enumerate}[label=\alph*.]
\item Combination
\item Combustion
\item Decomposition
\item Exchange
\end{enumerate}
\end{minipage}\\[0.2cm]
\hline
\end{tabular}

\noindent\begin{tabular}{|p{0.9\linewidth}|}
\hline
\textbf{(19)} What term describes the process by which a solid dissolves until no more can dissolve under given conditions?\\[0.2cm]
\begin{minipage}[t]{0.85\linewidth}
\begin{enumerate}[label=\alph*.]
\item Crystallization
\item Saturation
\item Dilution
\item Precipitation
\end{enumerate}
\end{minipage}\\[0.2cm]
\hline
\end{tabular}

\noindent\begin{tabular}{|p{0.9\linewidth}|}
\hline
\textbf{(20)} A cylinder contains 0.850~mol of gas at 2.50~atm. What is the temperature (K) if the volume is 15.0~L?\\[0.2cm]

\begin{minipage}[t]{0.85\linewidth}
\begin{enumerate}[label=\alph*.]
\item 210.~K
\item 298~K
\item 425~K
\item 537~K
\end{enumerate}
\end{minipage}\\[0.2cm]
\hline
\end{tabular}

\noindent\begin{tabular}{|p{0.9\linewidth}|}
\hline
\textbf{(21)} A student notices that hexane and water do not mix when shaken together. Which type of intermolecular forces explains this behavior?\\[0.2cm]
\begin{minipage}[t]{0.85\linewidth}
\begin{enumerate}[label=\alph*.]
\item London dispersion forces dominate in hexane, but water is polar.
\item Both substances form hydrogen bonds with each other.
\item Hexane is polar, and water is nonpolar.
\item Dipole-dipole interactions attract the two liquids together.
\end{enumerate}
\end{minipage}\\[0.2cm]
\hline
\end{tabular}

\noindent\begin{tabular}{|p{0.9\linewidth}|}
\hline
\textbf{(22)} Which of the following represents a redox reaction?\\[0.2cm]
\begin{minipage}[t]{0.85\linewidth}
\begin{enumerate}[label=\alph*.]
\item HCl + NaOH $\rightarrow$ NaCl + H$_2$O\\[0.1cm]
\item Zn + 2HCl $\rightarrow$ ZnCl$_2$ + H$_2$\\[0.1cm]
\item CaCO$_3$ $\rightarrow$ CaO + CO$_2$\\[0.1cm]
\item NaOH + HNO$_3$ $\rightarrow$ NaNO$_3$ + H$_2$O\\[0.1cm]
\end{enumerate}

\end{minipage}\\[0.2cm]
\hline
\end{tabular}

\noindent\begin{tabular}{|p{0.9\linewidth}|}
\hline
\textbf{(23)} A solution has [H$_3$O$^+$] = $2.22\times10^{-6}$ M. Calculate the pH, pOH, and [OH$^-$]. \\[4cm]
\hline
\end{tabular}

\noindent\begin{tabular}{|p{0.9\linewidth}|}
\hline
\textbf{(24)} A buffer is prepared using 0.250~M HCO$_3^-$ and 0.400~M CO$_3^{2-}$ (p$K_a = 10.33$). What are the buffer pH and corresponding pOH?\\[0.2cm]
\begin{minipage}[t]{0.85\linewidth}
\begin{enumerate}[label=\alph*.]
\item pH = 10.13,\ pOH = 3.87
\item pH = 10.33,\ pOH = 3.67
\item pH = 10.73,\ pOH = 3.27
\item pH = 10.15,\ pOH = 3.85
\end{enumerate}
\end{minipage}\\[0.2cm]
\hline
\end{tabular}

\noindent\begin{tabular}{|p{0.9\linewidth}|}
\hline
\textbf{(25)} A 425~mL cylinder contains nitrogen at $2.50\times10^{2}$~kPa and 22.0$^\circ$C. What pressure (in atm) will result if the gas is moved to a 1.00~L container and heated to 125$^\circ$C, assuming moles constant?\\[0.2cm]

\begin{minipage}[t]{0.85\linewidth}
\begin{enumerate}[label=\alph*.]
\item 0.564~atm
\item 1.41~atm
\item 0.282~atm
\item 2.24~atm
\end{enumerate}
\end{minipage}\\[0.2cm]
\hline
\end{tabular}

\noindent\begin{tabular}{|p{0.9\linewidth}|}
\hline
\textbf{(26)} What type of reaction occurs when potassium iodide reacts with lead(II) nitrate to form lead(II) iodide and potassium nitrate?\\[0.2cm]
2KI + Pb(NO$_3$)$_2$ $\rightarrow$ 2KNO$_3$ + PbI$_2$\\[0.2cm]
\begin{minipage}[t]{0.85\linewidth}
\begin{enumerate}[label=\alph*.]
\item Exchange
\item Combination
\item Decomposition
\item Combustion
\end{enumerate}
\end{minipage}\\[0.2cm]
\hline
\end{tabular}

\noindent\begin{tabular}{|p{0.9\linewidth}|}
\hline
\textbf{(27)} When salt dissolves in water, why is this considered a physical change? \\
\begin{minipage}[t]{0.85\linewidth}
\begin{enumerate}[label=\alph*.]
    \item The process creates a new chemical substance with different bonds.
    \item The salt changes state but its chemical composition remains the same.
    \item The salt reacts with water to form a new compound.
    \item The ions rearrange permanently to form a new molecule.
\end{enumerate}
\end{minipage}\\[0.2cm]
\hline
\end{tabular}

\noindent\begin{tabular}{|p{0.9\linewidth}|}
\hline
\textbf{(28)} A buffer is made by mixing 0.200~M HF and 0.350~M NaF (p$K_a=3.17$). After preparing the buffer, what are the pH and resulting $[\mathrm{H_3O^+}]$?\\[0.2cm]
\begin{minipage}[t]{0.85\linewidth}
\begin{enumerate}[label=\alph*.]
\item pH = 2.93,\ $[\mathrm{H_3O^+}] = 1.18\times10^{-3}$~M
\item pH = 3.33,\ $[\mathrm{H_3O^+}] = 4.70\times10^{-4}$~M
\item pH = 3.17,\ $[\mathrm{H_3O^+}] = 6.80\times10^{-4}$~M
\item pH = 3.60,\ $[\mathrm{H_3O^+}] = 2.50\times10^{-4}$~M
\end{enumerate}
\end{minipage}\\[0.2cm]
\hline
\end{tabular}

\noindent\begin{tabular}{|p{0.9\linewidth}|}
\hline
\textbf{(29)} A solution has [$H_3O^+$] = $3.20\times10^{-6}$ M. Its pH is closest to: \\[0.2cm]
\begin{minipage}{0.85\linewidth}
\begin{enumerate}[label=\alph*.]
\item 4.49 \item 5.49 \item 6.49 \item 7.49
\end{enumerate}
\end{minipage}\\[0.2cm]
\hline
\end{tabular}

\noindent\begin{tabular}{|p{0.9\linewidth}|}
\hline
\textbf{(30)} A weak acid has Ka = $3.00\times10^{-7}$. Which pKa is closest? \\[0.2cm]
\begin{minipage}{0.85\linewidth}
\begin{enumerate}[label=\alph*.]
\item 3.96 \item 4.22 \item 5.44 \item 6.80
\end{enumerate}
\end{minipage}\\[0.2cm]
\hline
\end{tabular}

\noindent\begin{tabular}{|p{0.9\linewidth}|}
\hline
\textbf{(31)} Two rigid containers are connected by a closed valve. One container holds helium at 2.0~atm in 4.0~L, and the other holds neon at 4.0~atm in 2.0~L. When the valve is opened and gases mix, what is the total pressure in the combined 6.0~L system (assume constant temperature)?\\[0.2cm]
\begin{minipage}[t]{0.85\linewidth}
\begin{enumerate}[label=\alph*.]
\item 2.00~atm
\item 2.30~atm
\item 2.67~atm
\item 6.00~atm
\end{enumerate}
\end{minipage}\\[0.2cm]
\hline
\end{tabular}

\noindent\begin{tabular}{|p{0.9\linewidth}|}
\hline
\textbf{(32)} A solution with pH = 2.00 has a hydronium concentration that is how many times larger than a solution with pH = 4.00? \\[0.2cm]
\begin{minipage}{0.85\linewidth}
\begin{enumerate}[label=\alph*.]
\item 2x \item 10x \item 100x \item 1000x
\end{enumerate}
\end{minipage}\\[0.2cm]
\hline
\end{tabular}

\noindent\begin{tabular}{|p{0.9\linewidth}|}
\hline
\textbf{(33)} Which of the following reactions is a redox reaction?\\[0.2cm]
\begin{minipage}[t]{0.85\linewidth}
\begin{enumerate}[label=\alph*.]
\item Zn + CuSO$_4$ $\rightarrow$ ZnSO$_4$ + Cu\\[0.1cm]
\item AgNO$_3$ + NaCl $\rightarrow$ AgCl + NaNO$_3$\\[0.1cm]
\item CaCO$_3$ $\rightarrow$ CaO + CO$_2$\\[0.1cm]
\item H$_2$O + SO$_3$ $\rightarrow$ H$_2$SO$_4$\\[0.1cm]
\end{enumerate}
\end{minipage}\\[0.2cm]
\hline
\end{tabular}

\noindent\begin{tabular}{|p{0.9\linewidth}|}
\hline
\textbf{(34)} Which factor determines how much the freezing point of a solution is lowered?\\[0.2cm]
\begin{minipage}[t]{0.85\linewidth}
\begin{enumerate}[label=\alph*.]
\item The number of solute particles
\item The solute's color
\item The solvent's taste
\item The solute's shape
\end{enumerate}
\end{minipage}\\[0.2cm]
\hline
\end{tabular}

\noindent\begin{tabular}{|p{0.9\linewidth}|}
\hline
\textbf{(35)} What happens to the molarity of a solution when it is diluted with more solvent?\\[0.2cm]
\begin{minipage}[t]{0.85\linewidth}
\begin{enumerate}[label=\alph*.]
\item It increases
\item It decreases
\item It stays the same
\item It doubles
\end{enumerate}
\end{minipage}\\[0.2cm]
\hline
\end{tabular}

\noindent\begin{tabular}{|p{0.9\linewidth}|}
\hline
\textbf{(36)} What is the molarity of a solution containing 5.00~mol of KBr in 2.00~L of solution?\\[0.2cm]
\begin{minipage}[t]{0.85\linewidth}
\begin{enumerate}[label=\alph*.]
\item 2.00~M
\item 2.50~M
\item 5.00~M
\item 10.0~M
\end{enumerate}
\end{minipage}\\[0.2cm]
\hline
\end{tabular}

\noindent\begin{tabular}{|p{0.9\linewidth}|}
\hline
\textbf{(37)} What is the percent yield if 10.0~g of product are expected but only 8.20~g are obtained?\\[0.2cm]

\begin{minipage}[t]{0.85\linewidth}
\begin{enumerate}[label=\alph*.]
\item 82.0\%
\item 92.0\%
\item 80.0\%
\item 88.0\%
\end{enumerate}
\end{minipage}\\[0.2cm]
\hline
\end{tabular}

\noindent\begin{tabular}{|p{0.9\linewidth}|}
\hline
\textbf{(38)} Which species is the conjugate acid of NH$_3$? \\[0.2cm]
\begin{minipage}{0.85\linewidth}
\begin{enumerate}[label=\alph*.]
\item NH$_2^-$ \item NH$_4^+$ \item N$^{3-}$ \item NH$_2$OH
\end{enumerate}
\end{minipage}\\[0.2cm]
\hline
\end{tabular}

\noindent\begin{tabular}{|p{0.9\linewidth}|}
\hline
\textbf{(39)} Osmosis is the movement of solvent molecules from:\\[0.2cm]
\begin{minipage}[t]{0.85\linewidth}
\begin{enumerate}[label=\alph*.]
\item Low solute concentration to high solute concentration
\item High solute concentration to low solute concentration
\item High pressure to low pressure
\item High temperature to low temperature
\end{enumerate}
\end{minipage}\\[0.2cm]
\hline
\end{tabular}

\noindent\begin{tabular}{|p{0.9\linewidth}|}
\hline
\textbf{(40)} A balloon is filled with a gas mixture containing 0.6~mol O$_2$ and 0.4~mol N$_2$ at a total pressure of 1.0~atm. What is the partial pressure of O$_2$?\\[0.2cm]
\begin{minipage}[t]{0.85\linewidth}
\begin{enumerate}[label=\alph*.]
\item 0.4~atm
\item 0.6~atm
\item 1.0~atm
\item 1.5~atm
\end{enumerate}
\end{minipage}\\[0.2cm]
\hline
\end{tabular}

\noindent\begin{tabular}{|p{0.9\linewidth}|}
\hline
\textbf{(41)} How many grams of KNO$_3$ are required to prepare 250.0~mL of 0.200~M KNO$_3$ solution?\\[0.2cm]
\begin{minipage}[t]{0.85\linewidth}
\begin{enumerate}[label=\alph*.]
\item 2.53~g
\item 5.06~g
\item 10.1~g
\item 20.2~g
\end{enumerate}
\end{minipage}\\[0.2cm]
\hline
\end{tabular}

\noindent\begin{tabular}{|p{0.9\linewidth}|}
\hline
\textbf{(42)} Convert 250.0 cm/s to mi/h. (Use $1~\text{mi}=1.609~\text{km}$) \\
\begin{minipage}[t]{0.85\linewidth}
\begin{enumerate}[label=\alph*.]
    \item $5.59~\text{mi/h}$
    \item $15.6~\text{mi/h}$
    \item $0.155~\text{mi/h}$
    \item $1.56~\text{mi/h}$
\end{enumerate}
\end{minipage}\\[0.2cm]
\hline
\end{tabular}

\noindent\begin{tabular}{|p{0.9\linewidth}|}
\hline
\textbf{(43)} For the reaction $\text{R} + \text{S} \rightleftharpoons \text{RS}$, which change could temporarily increase Q?\\[0.2cm]
\begin{minipage}[t]{0.85\linewidth}
\begin{enumerate}[label=\alph*.]
\item Adding RS
\item Removing R
\item Removing S
\item Diluting all species equally
\end{enumerate}
\end{minipage}\\[0.2cm]
\hline
\end{tabular}

\noindent\begin{tabular}{|p{0.9\linewidth}|}
\hline
\textbf{(44)} In the reaction $\text{HCl} + \text{H}_2\text{O} \rightarrow \text{H}_3\text{O}^+ + \text{Cl}^-$, which is the Brönsted-Lowry acid?\\[0.2cm]
\begin{minipage}[t]{0.85\linewidth}
\begin{enumerate}[label=\alph*.]
\item HCl
\item H$_2$O
\item H$_3$O$^+$
\item Cl$^-$
\end{enumerate}
\end{minipage}\\[0.2cm]
\hline
\end{tabular}

\noindent\begin{tabular}{|p{0.9\linewidth}|}
\hline
\textbf{(45)} Which of the following changes how much a solution’s freezing point is lowered?\\[0.2cm]
\begin{minipage}[t]{0.85\linewidth}
\begin{enumerate}[label=\alph*.]
\item How much solute is added
\item The type of solute (number of particles it makes)
\item A \& B
\item The color of the solute
\end{enumerate}
\end{minipage}\\[0.2cm]
\hline
\end{tabular}

\noindent\begin{tabular}{|p{0.9\linewidth}|}
\hline
\textbf{(46)} Convert $745~\text{torr}$ to atm. (Use $1~\text{atm}=760~\text{torr}$) \\
\begin{minipage}[t]{0.85\linewidth}
\begin{enumerate}[label=\alph*.]
    \item $0.980~\text{atm}$
    \item $1.02~\text{atm}$
    \item $0.760~\text{atm}$
    \item $0.745~\text{atm}$
\end{enumerate}
\end{minipage}\\[0.2cm]
\hline
\end{tabular}

\noindent\begin{tabular}{|p{0.9\linewidth}|}
\hline
\textbf{(47)} A 25\% (v/v) ethanol solution means:\\[0.2cm]
\begin{minipage}[t]{0.85\linewidth}
\begin{enumerate}[label=\alph*.]
\item 25~mL ethanol in 100~mL solution
\item 25~mL ethanol in 100~mL solvent
\item 25~g ethanol in 100~mL solution
\item 25~mL solution in 100~mL ethanol
\end{enumerate}
\end{minipage}\\[0.2cm]
\hline
\end{tabular}

\noindent\begin{tabular}{|p{0.9\linewidth}|}
\hline
\textbf{(48)} Which species is amphiprotic? \\[0.2cm]
\begin{minipage}{0.85\linewidth}
\begin{enumerate}[label=\alph*.]
\item OH$^-$ \item HSO$_4^-$ \item Cl$^-$ \item NH$_4^+$
\end{enumerate}
\end{minipage}\\[0.2cm]
\hline
\end{tabular}

\noindent\begin{tabular}{|p{0.9\linewidth}|}
\hline
\textbf{(49)} A buffer is prepared by mixing 0.300~mol acetic acid (p$K_a=4.76$) with 0.450~mol sodium acetate in 1.00~L solution. What is the pH of the buffer, and what is the corresponding pOH?\\[0.2cm]
\begin{minipage}[t]{0.85\linewidth}
\begin{enumerate}[label=\alph*.]
\item pH = 4.94,\ pOH = 9.06
\item pH = 4.76,\ pOH = 9.24
\item pH = 4.58,\ pOH = 9.42
\item pH = 5.12,\ pOH = 8.88
\end{enumerate}
\end{minipage}\\[0.2cm]
\hline
\end{tabular}

\noindent\begin{tabular}{|p{0.9\linewidth}|}
\hline
\textbf{(50)} For the equilibrium $ \text{PCl}_5(g) \rightleftharpoons \text{PCl}_3(g) + \text{Cl}_2(g)$, increasing pressure causes the reaction to shift: \\[0.2cm]
\begin{minipage}[t]{0.85\linewidth}
\begin{enumerate}[label=\alph*.]
\item Left, toward PCl$_5$
\item Right, toward PCl$_3$ and Cl$_2$
\item No shift because gases are unaffected by pressure
\item Toward whichever side has the higher total mass
\end{enumerate}
\end{minipage}\\[0.2cm]
\hline
\end{tabular}

\noindent\begin{tabular}{|p{0.9\linewidth}|}
\hline
\textbf{(51)} For each of the following, select the option that correctly identifies each property as physical (P) or chemical (C): \
\begin{minipage}[t]{0.85\linewidth}
\begin{enumerate}[label=\alph*.]
\item melting point of ice
\item flammability of hydrogen gas
\item density of aluminum
\item ability of iron to rust
\end{enumerate}
\textbf{Choices:} \
A) P, C, P, C \
B) C, P, C, P \
C) P, P, C, C \
D) C, C, P, P
\end{minipage}\\[0.2cm]
\hline
\end{tabular}

\noindent\begin{tabular}{|p{0.9\linewidth}|}
\hline
\textbf{(52)} For each of the following, select the option that correctly classifies each as a pure substance (PS) or a mixture (M): \
\begin{minipage}[t]{0.85\linewidth}
\begin{enumerate}[label=\alph*.]
\item sodium chloride
\item air
\item copper
\item seawater
\end{enumerate}
\textbf{Choices:} \
A) PS, M, PS, M \
B) M, PS, M, PS \
C) PS, PS, M, M \
D) M, M, PS, PS
\end{minipage}\\[0.2cm]
\hline
\end{tabular}

\noindent\begin{tabular}{|p{0.9\linewidth}|}
\hline
\textbf{(53)} A gas occupies 500.~mL at 25$^\circ$C. What volume will it occupy at 100.$^\circ$C if the pressure remains constant? \\[0.2cm]

\begin{minipage}[t]{0.85\linewidth}
\begin{enumerate}[label=\alph*.]
\item 583~mL
\item 626~mL
\item 665~mL
\item 740.~mL
\end{enumerate}
\end{minipage}\\[0.2cm]
\hline
\end{tabular}

\noindent\begin{tabular}{|p{0.9\linewidth}|}
\hline
\textbf{(54)} Convert 5850~$\text{cm}^2$ to $\text{m}^2$. \\
\begin{minipage}[t]{0.85\linewidth}
\begin{enumerate}[label=\alph*.]
    \item $0.0585~\text{m}^2$
    \item $5.85~\text{m}^2$
    \item $585~\text{m}^2$
    \item $0.000585~\text{m}^2$
\end{enumerate}
\end{minipage}\\[0.2cm]
\hline
\end{tabular}

\noindent\begin{tabular}{|p{0.9\linewidth}|}
\hline
\textbf{(55)} Which statement correctly defines ionization for a base in water? \\[0.2cm]
\begin{minipage}[t]{0.85\linewidth}
\begin{enumerate}[label=\alph*.]
\item Forming OH$^-$ by reacting chemically with water rather than releasing pre-existing ions
\item Separating into pre-existing ions without chemical change
\item Producing H$_3$O$^+$ by transferring a proton to water
\item Splitting into neutral molecules without forming ions
\end{enumerate}
\end{minipage}\\[0.2cm]
\hline
\end{tabular}

\noindent\begin{tabular}{|p{0.9\linewidth}|}
\hline
\textbf{(56)} Which term describes an acid that ionizes partially in water? \\[0.2cm]
\begin{minipage}[t]{0.85\linewidth}
\begin{enumerate}[label=\alph*.]
\item Strong acid \item Weak acid \item Strong base \item Neutral compound
\end{enumerate}
\end{minipage}\\[0.2cm]
\hline
\end{tabular}

\noindent\begin{tabular}{|p{0.9\linewidth}|}
\hline
\textbf{(57)} Osmotic pressure depends on which of the following factors?\\[0.2cm]
\begin{minipage}[t]{0.85\linewidth}
\begin{enumerate}[label=\alph*.]
\item The number of solute particles in solution
\item The color of the solute
\item The shape of the solvent molecules
\item The size of the container
\end{enumerate}
\end{minipage}\\[0.2cm]
\hline
\end{tabular}

\noindent\begin{tabular}{|p{0.9\linewidth}|}
\hline
\textbf{(58)} A buffer is formed by mixing 0.500~mol formic acid (p$K_a=3.75$) with 0.300~mol sodium formate. What are the buffer pH and the resulting $[\mathrm{H_3O^+}]$?\\[0.2cm]
\begin{minipage}[t]{0.85\linewidth}
\begin{enumerate}[label=\alph*.]
\item pH = 3.52,\ $[\mathrm{H_3O^+}] = 3.0\times10^{-4}$~M
\item pH = 3.75,\ $[\mathrm{H_3O^+}] = 1.8\times10^{-4}$~M
\item pH = 3.32,\ $[\mathrm{H_3O^+}] = 4.8\times10^{-4}$~M
\item pH = 3.90,\ $[\mathrm{H_3O^+}] = 1.3\times10^{-4}$~M
\end{enumerate}
\end{minipage}\\[0.2cm]
\hline
\end{tabular}

\noindent\begin{tabular}{|p{0.9\linewidth}|}
\hline
\textbf{(59)} A tectonic plate moves 860 km per day. What is its speed in m/s? \\
\begin{minipage}[t]{0.85\linewidth}
\begin{enumerate}[label=\alph*.]
    \item $9.95~\text{m/s}$
    \item $0.0995~\text{m/s}$
    \item $9.95 \times 10^3~\text{m/s}$
    \item $2.47~\text{m/s}$
\end{enumerate}
\end{minipage}\\[0.2cm]
\hline
\end{tabular}

\noindent\begin{tabular}{|p{0.9\linewidth}|}
\hline
\textbf{(60)} Which salt will produce a basic solution upon dissolution due to its parent species? \\[0.2cm]
\begin{minipage}{0.85\linewidth}
\begin{enumerate}[label=\alph*.]
\item NaF \item NH$_4$Br \item KCl \item Ca(NO$_3$)$_2$
\end{enumerate}
\end{minipage}\\[0.2cm]
\hline
\end{tabular}
\end{document}



